\subsection*{Step Further: \textbf{Symbolic Dimension}}

In Croqtile, \textbf{symbolic dimensions} give names to dimensions whose values are unknown at compile time. For example:

\begin{lstlisting}[language=croqtile,caption={Step Further: Symbolic Dimension}]
__co__ auto Matmul(s32 [M, K] lhs, s32 [K, N] rhs) {
  int tile_m = 32, tile_m = 8, tile_k = 16;
  s32 [M / tile_m, K / tile_k] tiled_lhs;
  s32 [K / tile_k, N / tile_n] tiled_rhs;
  // ...
}
\end{lstlisting}

Each input dimension has a symbolic name: \texttt{M}, \texttt{K}, or \texttt{N}. Unlike anonymous dimensions, these names make relationships between shapes explicit, improving readability and maintenance. They also give the compiler more information for consistency checks.
